\documentclass[a4paper]{article}

% Basic packages
% Español (México) con XeLaTeX: fontspec para fuentes Unicode, polyglossia
% para la división silábica y los nombres del documento en la variante mexicana.
\usepackage{fontspec}
\usepackage{polyglossia}
\setdefaultlanguage[variant=mexican]{spanish}
% Los nombres chinos van como «pinyin (original)»: xeCJK da a los caracteres
% Han su propia fuente; sin ella XeLaTeX los omite con un aviso.
% FandolSong no cubre algunos caracteres de nombres (U+73FA); WenQuanYi Zen Hei sí.
\usepackage[AutoFallBack=true]{xeCJK}
\setCJKmainfont{FandolSong-Regular.otf}
\setCJKfallbackfamilyfont{\CJKrmdefault}{WenQuanYi Zen Hei}
% polyglossia no traduce los nombres de listings.
\AtBeginDocument{\providecommand{\lstlistingname}{}\renewcommand{\lstlistingname}{Listado}%
  \providecommand{\lstlistlistingname}{}\renewcommand{\lstlistlistingname}{Índice de listados}}
\usepackage{amsmath, amssymb}
\usepackage{graphicx}
\usepackage[margin=2.5cm]{geometry}
\usepackage[most]{tcolorbox}
\usepackage{etoolbox}
\usepackage{listings}
\usepackage{hyperref}
\usepackage{booktabs}
\usepackage{subcaption}
\usepackage{float}
\usepackage{tikz}
\IfFileExists{pgfplots.sty}{
    \usepackage{pgfplots}
    \pgfplotsset{compat=1.18}
}{}

% Highlight boxes
\newtcolorbox{knowledgebox}[1]{
    enhanced, colback=blue!5!white, colframe=blue!75!black, colbacktitle=blue!75!black,
    coltitle=white, fonttitle=\bfseries, title={#1}, attach boxed title to top left={yshift=-2mm, xshift=2mm},
    boxrule=1pt, sharp corners
}

\newtcolorbox{importantbox}[1]{
    enhanced, colback=yellow!10!white, colframe=yellow!80!black, colbacktitle=yellow!80!black,
    coltitle=black, fonttitle=\bfseries, title={#1}, sharp corners
}

\newtcolorbox{warningbox}[1]{
    enhanced, colback=red!5!white, colframe=red!75!black, colbacktitle=red!75!black,
    coltitle=white, fonttitle=\bfseries, title={#1}, sharp corners
}

% Listings
\lstset{
    language=Python,
    basicstyle=\ttfamily\small,
    keywordstyle=\color{blue},
    stringstyle=\color{red!60!black},
    commentstyle=\color{green!60!black},
    breaklines=true,
    frame=single,
    numbers=left,
    numberstyle=\tiny\color{gray},
    captionpos=b,
    extendedchars=true
}

% Front-page metadata
\newcommand{\notetitle}{Notas de entrevista: Hu Yuanming (胡渊鸣) — Meshy AI, taichí, MIT, la Clase Yao de Tsinghua, gráficos por computadora y filosofía del emprendimiento}
\newcommand{\noteauthors}{Elaboradas a partir de materiales públicos del curso}
\newcommand{\notedate}{2025-08-08}
\newcommand{\videochannel}{WhynotTV}
\newcommand{\videopublishdate}{2025-08-08}
\newcommand{\videoduration}{184 minutos}
\newcommand{\videourl}{https://www.bilibili.com/video/BV1XmtyzKEzQ/}
\newcommand{\videocoverpath}{cover.jpg}
\newcommand{\repourl}{https://github.com/hqhq1025/ai-course-notes}

% Footer with repo info
\usepackage{fancyhdr}
\pagestyle{fancy}
\fancyhf{}
\fancyfoot[C]{\thepage}
\fancyfoot[R]{\footnotesize\color{gray}\href{\repourl}{AI Course Notes}}
\renewcommand{\headrulewidth}{0pt}
\renewcommand{\footrulewidth}{0pt}

\begin{document}

\begin{titlepage}
\centering
{\Large Notas de entrevista\par}
\vspace{1.2cm}
{\huge\bfseries \notetitle\par}
\vspace{0.8cm}
{\large \noteauthors\par}
\vspace{0.3cm}
{\large \notedate\par}
\vspace{1.2cm}

\ifdefempty{\videocoverpath}{
{\small Al generar las notas, indique la ruta local de la portada del video.\par}
}{
\includegraphics[width=0.82\textwidth,height=0.45\textheight,keepaspectratio]{\videocoverpath}\par
}

\vfill
\begin{tcolorbox}[width=0.9\textwidth, colback=black!2!white, colframe=black!60, sharp corners]
\textbf{Autor o canal del video}: \videochannel\par
\textbf{Fecha de publicación}: \videopublishdate\par
\textbf{Duración del video}: \videoduration\par
\ifdefempty{\videourl}{}{
\textbf{Enlace del video}: \href{\videourl}{\nolinkurl{\videourl}}\par
}\par
\textbf{Página del proyecto}: \href{\repourl}{\nolinkurl{\repourl}}\par
\end{tcolorbox}
\end{titlepage}

\tableofcontents
\newpage

%% --- Inicio del contenido --- %%

\section{Introducción: quién es Hu Yuanming (胡渊鸣)}

Este episodio es el segundo de WhynotTV Podcast: el conductor Tairan (泰然) sostuvo una conversación a fondo de tres horas con Hu Yuanming (胡渊鸣), CEO de Meshy AI. La trayectoria de Hu Yuanming es sumamente rica: escribió programas de simulación física desde la primaria, cursó la licenciatura en la Clase Yao (姚班) de Tsinghua, durante su doctorado en el MIT creó de manera independiente el lenguaje de programación Taichi (太极), tras doctorarse fundó Taichi Graphics (太极图形), y tras 18 meses de exploración de la comercialización dio un giro hacia la plataforma de IA generativa 3D Meshy, que hoy atiende a más de 4 millones de usuarios, con un crecimiento anual de ingresos superior a 10 veces.

Cuando le preguntaron «quién es Hu Yuanming», su respuesta fue breve y profunda:

\begin{importantbox}{Autodefinición}
«Soy una persona que actúa impulsada por su propio interés y su sentido de misión.» Esta frase atraviesa toda la entrevista: desde escribir juegos en la primaria hasta crear Taichi y después fundar Meshy, el impulso del interés ha sido siempre la lógica central de las acciones de Hu Yuanming.
\end{importantbox}

El contenido de la entrevista abarca: la reflexión sobre la naturaleza de la tecnología y del mundo, la autosubversión durante el doctorado, el giro de vida o muerte al emprender, y la filosofía evolutiva de «matar cada tres meses a la versión anterior de sí mismo». A continuación se organizan por tema las ideas centrales y las lecciones que ahí aparecen.

\begin{table}[H]
\centering
\caption{Cronología de la vida de Hu Yuanming}
\begin{tabular}{lll}
\toprule
\textbf{Periodo} & \textbf{Evento} & \textbf{Palabra clave} \\
\midrule
Preescolar/primaria & Escribió simulaciones físicas y juegos con VB6.0 & Germen del interés \\
Secundaria & Competencias de programación (NOI), sistemas de masas y resortes & Acumulación de bases \\
Preparatoria & Simulación de cuerpos rígidos, fichas de dominó & Descubrió los límites del modelado \\
2013--2017 & Licenciatura en la Clase Yao de Tsinghua & Prácticas en MSRA, artículo en CVPR \\
2017--2021 & Doctorado en el MIT (CSAIL) & Taichi, DiffTaichi \\
2020 & Prácticas en NVIDIA, curso Games201 & Práctica de compiladores \\
2021 & Se doctoró, fundó Taichi Graphics & Comenzó a emprender \\
2021--2022 & Exploración de la comercialización de Taichi & 18 meses de búsqueda \\
Segunda mitad de 2022 & Vio ChatGPT/SD, Pivot hacia Meshy & Primeros principios \\
2023 & Iteración rápida del producto Meshy & Lanzamiento en 8 horas \\
2024--2025 & Meshy con 4 millones de usuarios, crecimiento anual superior a 10 veces & Primer lugar del mercado \\
\bottomrule
\end{tabular}
\end{table}
\section{Crecimiento y formación: de simular el mundo a entenderlo}

\subsection{Infancia: crear un mundo virtual con suma, resta, multiplicación y división}

Los padres de Hu Yuanming (胡渊鸣) se dedicaban a la enseñanza de informática, así que él tuvo contacto con las computadoras desde el jardín de niños y jugó títulos como Xianjian Qixia Zhuan, Red Alert y Age of Empires. Estas experiencias dieron origen a una intuición central: \textbf{con suma, resta, multiplicación y división se puede crear un mundo dentro de una computadora}.

Una reflexión filosófica más profunda ya germinaba en esa época. Recuerda que, de niño, observaba el melocotonero de abajo y se preguntaba: "Cuando no estoy observando este melocotonero, ¿sigue creciendo en silencio, o se hace el flojo y permanece igual, y solo adopta la forma que debería tener en el instante en que lo miro?" Más tarde comprendió que, al ser parte de lo simulado, uno no puede distinguir entre ambos casos: algo análogo al «efecto del observador» de la mecánica cuántica.

\begin{knowledgebox}{Hipótesis de la simulación (Simulation Hypothesis)}
Hu Yuanming extendió esa intuición filosófica infantil a su comprensión del mundo físico: la constante de Planck podría ser el límite de precisión de punto flotante de un sistema simulado, la velocidad de la luz podría ser un límite de distancia de interacción impuesto por recursos computacionales insuficientes, y el Big Bang podría no ser más que la configuración de las condiciones iniciales. Esta forma de pensar influyó de manera profunda en su visión posterior de la simulación física y la IA.
\end{knowledgebox}

Al entrar a la primaria, escribió código de simulación física en Visual Basic 6.0 y, sin haber estudiado todavía las leyes de Newton, "descubrió por su cuenta" la segunda ley de Newton mediante la programación: cuando movía un objeto directamente con las flechas del teclado le parecía poco realista, así que le definió un incremento de velocidad al objeto, y solo después supo que eso era $F = ma$.

\begin{table}[H]
\centering
\caption{Experiencias de programación y simulación de Hu Yuanming en cada etapa}
\begin{tabular}{lll}
\toprule
\textbf{Etapa} & \textbf{Herramienta/lenguaje} & \textbf{Contenido de la simulación} \\
\midrule
Jardín de niños/primaria & Visual Basic 6.0 & Juego de nave espacial, simulación física básica \\
Primaria/secundaria & RPG Maker XP, Ruby & Juego de rol, sistema de masa-resorte (Mass-Spring) \\
Preparatoria & C++ & Simulación de cuerpo rígido (Rigid Body), fichas de dominó \\
Licenciatura (Tsinghua) & C++, CUDA & Simulación de fluidos, reproducción de artículos de SIGGRAPH \\
Doctorado (MIT) & Taichi, LLVM & MPM, simulación diferenciable, compilador \\
\bottomrule
\end{tabular}
\end{table}

En la secundaria construyó un sistema de masa-resorte (Mass-Spring System), inspirado en el juego World of Goo: se tienen muchas bolas que se pueden conectar entre sí para construir estructuras, y su representación en la computadora es precisamente un mass-spring system. En la preparatoria comenzó a estudiar la simulación de cuerpo rígido (Rigid Body Simulation) e intentó simular fichas de dominó en la computadora: "porque en el mundo real, si apilas algo, no tienes el espacio, y basta con tocarlo sin querer para que todo se venga abajo".

En ese proceso descubrió algo: \textbf{los modelos de colisión y fricción entre dos cuerpos rígidos son siempre aproximados}. Esta comprensión influyó después de manera profunda en cómo entendió las limitaciones de la simulación, y también sentó las bases de su trabajo posterior en simulación diferenciable durante el doctorado, así como de su postura abierta hacia la data-driven simulation.

Entre la secundaria y la preparatoria, buena parte de su tiempo lo ocupaban las competencias de informática (como la NOI); el competitive programming entrenó su capacidad de "escribir código correcto bajo mucha presión", una habilidad que "hasta hoy sigue beneficiándome".

\subsection{La influencia del competitive programming en emprender}

Aunque Hu Yuanming no menciona explícitamente la relación entre el competitive programming y emprender, en la entrevista pueden verse varias conexiones implícitas:

\begin{enumerate}
\item \textbf{Resolver problemas con tiempo limitado $\rightarrow$ iteración rápida}: las competencias exigen escribir código correcto en un tiempo limitado, y emprender exige probar y iterar rápido. Meshy pasó de cero a estar en línea en solo 8 horas.
\item \textbf{Pensamiento de casos límite $\rightarrow$ diseño de sistemas}: quienes compiten se acostumbran a considerar todo tipo de corner cases, un tipo de pensamiento sumamente importante al diseñar compiladores y sistemas distribuidos.
\item \textbf{Presión de la competencia $\rightarrow$ capacidad de resistir presión}: mantener la calma y la eficiencia en entornos de alta presión es una exigencia común tanto de las competencias como de emprender.
\item \textbf{Cultura de las clasificaciones $\rightarrow$ orientación a resultados}: las competencias dan una retroalimentación clara mediante clasificaciones, lo que cultiva el hábito de "respetar los hechos y no engañarse a uno mismo".
\end{enumerate}

\subsection{La clase Yao de Tsinghua: descubrir que siempre hay algo más allá}

Al entrar a la clase Yao (姚班) de Tsinghua, el mayor tesoro que obtuvo Hu Yuanming fue "conocer a un grupo de compañeros extraordinariamente capaces". Había quien apenas estudiaba y aun así sacaba 100 en los exámenes, y quien ganaba con facilidad la medalla de oro del ACM World Final. Compartió una experiencia de frustración que lo marcó:

\begin{quote}
"La primera materia se llamaba 'Matemáticas Aplicadas a la Computación'. Por lo general, al entrar a Tsinghua primero se estudia cálculo, álgebra lineal y probabilidad, pero esa materia parecía asumir directamente que ya sabías todo eso y empezaba de una vez. Lo más absurdo fue que el profesor dijo: 'Este tema no me dio tiempo de explicarlo en clase, así que lo puse como pregunta de examen, con la esperanza de que lo dominen para cuando lleguen al examen.' Después del examen todos decían que había estado dificilísimo. Un compañero dijo: 'voy a reprobar, estuvo dificilísimo, dificilísimo', y al final resultó que sacó 98.5."
\end{quote}

\begin{importantbox}{Conocer antes a gente extraordinaria}
"Conocer antes a gente extraordinaria le ayuda mucho al crecimiento de una persona. Te hace darte cuenta de que siempre hay algo más allá, de que siempre hay alguien mejor, de que uno no es tan capaz como se cree y de que todavía hay mucho margen para mejorar." Esta forma de entender las cosas no comenzó en la clase Yao, sino que ya la había vivido desde la secundaria, al competir y formar parte de la selección de la provincia de Jiangsu.
\end{importantbox}

En Tsinghua, las oportunidades para dedicarse a la gráfica por computadora en China eran escasas, y la mayoría de los compañeros más destacados se orientaban hacia la IA (justo cuando estallaban, uno tras otro, AlexNet, ResNet y AlphaGo). Hu Yuanming pasaba prácticamente todo su tiempo en el dormitorio leyendo artículos de SIGGRAPH y reproduciendo sus resultados; esta investigación independiente, de manera inesperada, entrenó la capacidad de "empezar desde cero" que más tarde usaría al emprender.

También hizo una estancia en Microsoft Research Asia (MSRA), donde, en el grupo de gráfica de redes del profesor Steve Lin, elaboró un artículo para CVPR (predecir el balance de blancos con redes neuronales) y un trabajo de posprocesamiento de imágenes con GAN+RL, este último entrenó más de 1,000 modelos para apenas obtener un resultado utilizable. Esa experiencia le dio una visión lúcida de la madurez que tenía entonces la tecnología de IA: el problema del sparse reward en RL "en realidad sigue sin resolverse hasta hoy"; "tampoco hoy: la semana pasada nos pasamos toda una semana ajustando a mano la reward function".

Un detalle digno de notarse: durante su licenciatura en Tsinghua, \textbf{todos} los artículos que logró publicar fueron de la línea de IA: "de Graphics no publiqué ni uno". Aunque lo que más le apasionaba era Graphics, la línea de IA le permitía obtener resultados con más facilidad. Esto también anticipaba la tendencia posterior: la fusión profunda entre Graphics e IA.

\begin{warningbox}{No confiar demasiado en los artículos}
"La investigación tiene muchas limitaciones, y no se puede confiar demasiado en lo que dicen los artículos." Al reproducir artículos de SIGGRAPH descubrió que muchos de ellos no explicaban con claridad todos los detalles, una lección que después lo llevó, en el proyecto Taichi, a insistir en que "todos los resultados experimentales deben poder reproducirse de extremo a extremo con un solo comando".
\end{warningbox}

\subsection{Doctorado en el MIT: de casi abandonar a crear Taichi}

En 2017 entró al MIT, y su primer año fue muy doloroso: sus intereses de investigación no coincidían con los de su asesor, y tenía que dedicar mucho tiempo a cosas que no disfrutaba. Se describe como alguien que "puede dar el 1,000\% cuando hace algo que disfruta, pero sufre mucho al hacer algo que no disfruta".

El punto de inflexión llegó gracias al consejo de dos compañeros mayores de Tsinghua: Jiajun le sugirió "convertir el simulator en differentiable", lo que más tarde dio origen a dos artículos importantes, ChainQueen y DiffTaichi; Junyan, por su parte, insistió en la importancia de "hacer research con impact". Al final cambió al grupo de Fredo Durand y Bill Freeman, dos asesores bastante hands-off que le dieron plena libertad de investigación.

\begin{importantbox}{Cuatro cambios que haría si repitiera el doctorado}
\begin{enumerate}
\item \textbf{Elegir el rumbo}: no ir a NVIDIA a escribir Taichi, sino aceptar la oferta de Kaiming He (何恺明) para ir a FAIR a hacer investigación en IA, y así tener contacto más temprano con el frontier de la IA
\item \textbf{Ampliar el alcance}: seguir trabajando en Taichi, pero sin limitarse a la physical simulation, y avanzar antes hacia la optimización de Transformer
\item \textbf{Meta ambiciosa}: seguir teniendo como meta graduarse en un año; "aspira a lo alto y llegarás a lo medio; aspira a lo medio y llegarás a lo bajo; aspira a lo bajo y no llegarás a nada"
\item \textbf{Más intercambio}: dedicar más tiempo a conversar con gente destacada de distintos campos, en lugar de encerrarse todos los días en el dormitorio a escribir Taichi
\end{enumerate}
\end{importantbox}

Sobre este último punto, citó una idea de Richard Hamming: "Si te encierras en tu oficina, cierras la puerta y haces tus propias cosas todos los días, en el corto plazo, claro, produces más. Pero cinco años después, sin duda no sabrás qué hacer." Solo dejando la puerta abierta y escuchando en qué están trabajando los demás se pueden descubrir los problemas importantes que vale la pena resolver.

\subsection{MIT y Stanford: cómo la ubicación geográfica influye en el rumbo académico}

Hu Yuanming tiene una observación singular sobre la diferencia entre el MIT y Stanford en la era de la IA. Cuando ingresó en 2017, el CSAIL del MIT no tenía a muchas personas trabajando en IA; el ambiente emprendedor de toda la región de Boston giraba sobre todo en torno a la Biotech, y faltaba un impulso desde la industria de la IA.

\begin{knowledgebox}{Cómo el ecosistema industrial moldea el rumbo académico}
"Comparadas con escuelas cercanas a un centro industrial como Stanford, escuelas como el MIT en realidad tienen una aceptación mucho menor de la IA. Al no haber un impulso desde la industria, en toda la región de Boston se hace más bien Bio. Las empresas que salen de ahí son básicamente del campo de la Bio." Esto hizo que el MIT no se mantuviera especialmente al día durante toda la era de la IA. "Esto puede ser algo bueno o algo malo; hay que verlo de manera dialéctica."
\end{knowledgebox}

Esta observación aporta una idea importante para entender el ecosistema académico: el rumbo de investigación de una universidad no depende solo de los intereses de sus profesores, sino que está profundamente influido por el ecosistema industrial que la rodea. Que Stanford se haya convertido en un bastión de la IA está estrechamente ligado a la industria tecnológica de Silicon Valley.

\subsection{La «regla del 95/5» para hacer investigación}

En un artículo de Zhihu, Hu Yuanming citó una frase de Einstein: "Si me dieran una hora para resolver un problema, pasaría 55 minutos pensando cuál es la pregunta correcta, y solo 5 minutos resolviéndola." Esta «regla del 95/5» se convirtió en la síntesis central de su metodología de investigación.

En cada entrevista de trabajo para Meshy pregunta: "¿lo que haces es importante?" Cuando a él mismo le preguntaron cuál era la cosa más importante que había hecho recientemente, respondió: \textbf{reclutamiento} (40\% del tiempo), \textbf{pensar en la estrategia futura} (30\%) y \textbf{participar en las operaciones cotidianas} (30\%).

\begin{warningbox}{La mayoría de los doctorados siguen "inflando la producción" en la etapa dos}
"Publicas un artículo, sientes que quedó muy bien, pero cuando hablas con otros no les resulta nada exciting; tu artículo solo le interesa a unas cuantas personas. ¿Por qué? Porque no has terminado de pensar cuál es el rumbo realmente importante. Solo después de suficientes fracasos y reflexión llegas a saber cómo debe ser un buen rumbo." Muchos doctorados, incluso después de tener tres o cuatro artículos en conferencias de primer nivel, siguen usando el pensamiento de la etapa uno, «acumular producción», para inflar publicaciones, sin cambiar al modo de la etapa dos: «buscar el problema importante».
\end{warningbox}

Pone a Kaiming (Kaiming He, 何恺明) como ejemplo de un «doctorado exitoso de etapa dos»: ResNet no es más que "un signo de más", pero su influencia supera con creces la de decenas de artículos comunes. En su artículo de Zhihu escribió: "no te conformes con ser un pez pequeño dentro de SIGGRAPH; ve a ser un pez grande en el estanque de todo el CS". Si Taichi 2.0 logra tener una aplicación amplia en la era de los modelos grandes y su influencia llega a superar 50 veces la de la versión 1.0, considera que "eso sí sería convertirse en un pez grande".

\subsection{Sobre la suerte: el 99\% de la atribución}

Hu Yuanming insiste una y otra vez en la importancia de la suerte. Cuando le preguntaron si el hecho de haber sabido tan temprano qué disfrutaba era talento o suerte, su respuesta invita a la reflexión:

\begin{importantbox}{La suerte como sistema operativo de base}
"Creo que el talento también es suerte. Simplemente te tocó, genetically, un gen bastante bueno, y resulta que esta época necesitaba justo lo que ese gen te permite hacer. Así que creo que tal vez el 99\% sea suerte. Incluso el que puedas esforzarte también es parte de la suerte, porque gracias a esa suerte sabes qué te gusta, y por eso puedes esforzarte."
\end{importantbox}

Esta forma tan marcada de atribuirlo todo a la suerte no es negativa; de hecho, es la máxima expresión de ser humble: reconocer lo azaroso de las propias ventajas y mantener empatía hacia quienes han tenido menos suerte. Esto también explica por qué no juzga a quienes eligen "tirarse a la bardas" (躺平): "es una elección personal, y también está bien si a esa persona le hace feliz."

\subsection{Resumen de la sección}

El camino de crecimiento de Hu Yuanming muestra varios rasgos clave: (1) descubrir su interés desde muy temprano y profundizar en él de manera constante; (2) "descubrir por su cuenta" las leyes físicas a través de la práctica, en lugar de aprenderlas de forma pasiva; (3) encontrarse en cada etapa con gente más capaz que él y mantenerse humble; (4) conservar un pensamiento crítico frente a la investigación académica. Su manera de manejar la relación con su asesor (del sufrimiento del primer año a la libertad tras cambiar de asesor) y su reflexión sobre las decisiones profesionales (arrepentirse de no haber ido a FAIR) muestran una capacidad genuina de autoexamen.
\section{Taichi (太极): un lenguaje de programación de una sola persona}

\subsection{Qué es la computación gráfica}

Hu Yuanming (胡渊鸣) define la computación gráfica de forma concisa y contundente: \textbf{crear un mundo virtual con sumas, restas, multiplicaciones y divisiones}. La CPU y la GPU solo pueden hacer sumas, restas, multiplicaciones, divisiones y funciones trigonométricas, pero con estos building blocks más básicos hay que crear imágenes de videojuego que parezcan vivas.

El desarrollo de la computación gráfica se divide, a grandes trazos, en varias etapas:

\begin{enumerate}
\item \textbf{Visualización y renderizado con realismo} (antes de aproximadamente 1995)
\item \textbf{Simulación física} (después de 1995): la simulación de humo de Fedkiw et al. (Visual Simulation of Smoke) abrió la era de la simulación
\item \textbf{Fotografía computacional}: técnicas como super resolution, indisociables de la visión por computadora
\item \textbf{Neural Graphics}: nuevos paradigmas impulsados por IA como DLSS, NeRF y Gaussian Splatting
\end{enumerate}

\begin{knowledgebox}{Señal para saber si un campo está en declive}
"Cuando notas que un problema se vuelve cada vez más limpio y cada vez más well-defined, eso indica que ya no queda mucho por hacer con ese problema. Cuando se forma un benchmark particularmente sistemático, cuando las soluciones se vuelven cada vez más complejas, eso indica que las oportunidades de abrir nuevo territorio ya son escasas." Hu Yuanming (胡渊鸣) ilustra esto con el ejemplo de AMCMCPPM (muestreo unificado de caminos por cadenas de Markov Monte Carlo adaptativas) en el campo del renderizado.
\end{knowledgebox}

\subsection{El nacimiento y la misión de Taichi}

La misión central de Taichi es construir una \textbf{infraestructura (Infrastructure)} que haga que la simulación sea más fácil de desarrollar, se ejecute más rápido, use menos memoria y se integre bien con el ecosistema circundante (Python, PyTorch).

Un logro emblemático fue ejecutar una simulación de mil millones ($10^9$) de partículas en una sola GPU 3090 (24 GB de memoria de video). Para lograrlo, el espacio de almacenamiento que ocupa cada partícula no puede superar los 24 bytes. Él hizo optimizaciones de cuantización extremas a nivel del compilador: las coordenadas XYZ de una partícula se representan con un entero de 32 bits (X con 11 bits, Y con 11 bits, Z con 10 bits); ese es el trabajo de QuantTaichi.

\begin{importantbox}{Por qué Taichi necesitaba construir su propio compilador}
El punto clave es que, aunque CUDA permite a los programadores de C++ escribir programas de GPU, la madurez del ecosistema de Python, de la infraestructura del compilador LLVM y de métodos de simulación como MPM creó la oportunidad de construir lenguajes específicos de dominio en un nivel más alto. Taichi se basa en la sintaxis de Python y, mediante un compilador construido desde cero, compila código de alto nivel a programas de GPU de alto rendimiento, de modo que el usuario no tiene que manejar manualmente operaciones de bits ni optimización de memoria.
\end{importantbox}

La experiencia de desarrollar Taichi de forma independiente en el MIT es descrita por Hu Yuanming (胡渊鸣) como "muy placentera": él desempeñaba a la vez cuatro roles:

\begin{table}[H]
\centering
\caption{Los cuatro roles en el desarrollo de Taichi}
\begin{tabular}{ll}
\toprule
\textbf{Rol} & \textbf{Trabajo concreto} \\
\midrule
Gerente de producto & Analizar las necesidades de los usuarios de simulación (porque él mismo era el usuario) \\
Científico & Conocer la tecnología de vanguardia y convertirla en características del producto \\
Ingeniero & CI/CD, arquitectura de software, optimización de rendimiento en C++, compilación multihilo, caché de compilación \\
CEO & Integrar los resultados de los otros tres roles para que más gente lo usara (promoción en Zhihu, etc.) \\
\bottomrule
\end{tabular}
\end{table}

"Usar mi propio lenguaje de programación para implementar mi propio programa y descubrir que además era más cómodo de usar que CUDA: creo que en ese momento me sentí un poco increíble." Este fue el momento que más lo emocionó durante el desarrollo de Taichi.

\subsection{Frozen en 99 líneas de código}

El demo más conocido de Taichi es el efecto de "Frozen en 99 líneas de código": con un código de Python extremadamente corto se implementó una simulación de nieve impulsada por el método de puntos materiales (MPM). Este demo obtuvo una atención enorme en redes sociales y se convirtió en una de las etiquetas más emblemáticas de Hu Yuanming (胡渊鸣).

Sin embargo, su actitud ante esta etiqueta es ambivalente: "Creo que hace tiempo dejaste de idealizar este tipo de etiquetas, y una etiqueta así de ninguna manera muestra quién eres en realidad." De hecho, el demo de 99 líneas es solo la punta del iceberg de lo que Taichi puede hacer: detrás hay toda una pila de compilador, un sistema de tipos, un motor de diferenciación automática, generadores de código para múltiples backends y mucha más infraestructura.

\begin{knowledgebox}{La programación diferenciable (Differentiable Programming) de Taichi}
Una característica clave de Taichi es que soporta la programación diferenciable, es decir, tu programa de simulación puede derivarse automáticamente. Esto significa que puedes:
\begin{itemize}
\item Dado un estado objetivo, optimizar de forma inversa los parámetros de la simulación
\item Incrustar el simulator en el ciclo de entrenamiento de una neural network
\item Usar gradient descent para optimizar la estrategia de control de un robot
\end{itemize}
Los dos artículos DiffTaichi y ChainQueen se basan justamente en esta característica para combinar simulation y robotics. Esta dirección fue sugerida por Jiajun (家俊), un compañero mayor de Tsinghua: "¿Por qué no intentas hacer que tu simulator sea differentiable?"
\end{knowledgebox}

Hu Yuanming (胡渊鸣) recuerda que basta con agregar una o dos líneas al código de Taichi para convertir una simulación ordinaria en una simulación diferenciable: "También pienso que solo construyendo tu propio compilador puedes lograr esto." Esta capacidad de abstracción a nivel de compilador es algo a lo que el código de CUDA escrito a mano no puede llegar.

\subsection{La visión de Taichi 2.0}

La entrevista reveló que la planeación de Taichi 2.0 está en marcha. El objetivo es que Taichi siga teniendo un uso amplio en la era de los modelos grandes, con una influencia que supere en más de 50 veces la de la versión 1.0. Las direcciones concretas podrían incluir:

\begin{itemize}
\item \textbf{Soporte para la optimización de Transformer}: si en su momento Taichi se hubiera orientado antes hacia esta dirección, "quizás ahora Taichi estaría en la posición de Triton"
\item \textbf{Hybrid Simulator}: combinar de verdad simulation y neural network, con una parte de simulación física y una parte data-driven
\item \textbf{Soporte para el ecosistema de LLM}: encontrar el escenario de aplicación de Taichi en el ámbito de los LLM
\end{itemize}

Un dato interesante: el número de GitHub Stars de Taichi supera al de proyectos como NVIDIA Warp, Triton y Mojo: "Taichi sigue siendo bastante popular, solo que no ha encontrado un ámbito grande donde desplegarse."

\subsection{Las tres evoluciones de la filosofía del código abierto}

La visión de Hu Yuanming (胡渊鸣) sobre el código abierto pasó por tres etapas:

\textbf{Primera etapa (durante el doctorado en el MIT)}: el código abierto es lo mejor. En ese entonces muchos artículos de Graphics no publicaban código abierto, y él insistió en que todos los resultados experimentales de Taichi pudieran reproducirse con un solo comando.

\textbf{Segunda etapa (inicio del emprendimiento)}: comercializar el código abierto es sumamente arriesgado. "Ya le diste a la gente lo mejor que tenías, entonces, ¿cómo vas a ganar dinero?" La inversión en Taichi disminuyó y el equipo se trasladó a trabajar en Meshy.

\textbf{Tercera etapa (actual)}: el código abierto es una herramienta estratégica de marketing.

\begin{importantbox}{El código abierto como mercadotecnia técnica}
"Es cierto que el código abierto no genera ganancias, pero un buen proyecto de código abierto puede ayudarte a atraer al mejor talento y a los mejores clientes: se convierte en un medio de marketing." En la industria de los videojuegos, la inversión en investigación y desarrollo y en promoción de mercado suele estar en una proporción de 1:1 o incluso mayor. Hacer mercadotecnia técnica a través del código abierto es el mejor camino para los fundadores técnicos que no son hábiles en el marketing tradicional. Pero esto requiere ser "muy strategic": hay que pensar en una buena forma de comercialización que lo complemente.
\end{importantbox}

Él reconoce con franqueza que el código abierto es una parte de su ideal como Founder de la que no puede desprenderse: "tarde o temprano es posible que vuelvas a tomar este camino del código abierto." Actualmente el equipo está planeando Taichi 2.0, posicionado como "infrastructure supporting Graphics and AI".

\subsection{Resumen de la sección}

El proyecto Taichi encarna el espíritu geek de Hu Yuanming (胡渊鸣) de "construir las cosas desde la capa más básica" y también es el puente que lo llevó de la academia al emprendimiento. En lo técnico, mostró el enorme potencial de un compilador específico de dominio (simulación de mil millones de partículas); en lo metodológico, forjó una capacidad integral de "gerente de producto + científico + ingeniero + CEO". Las tres evoluciones de la filosofía del código abierto reflejan, a su vez, el equilibrio entre el ideal y la realidad comercial.
\section{La fusión de la simulación física y la IA}

\subsection{La brecha entre los tres mundos}

En la entrevista hay un pasaje excelente sobre los «tres mundos»: (1) el mundo real; (2) el mundo que podemos observar a simple vista; (3) el mundo que modelamos con modelos físicos.

\begin{importantbox}{El Gap entre los tres mundos}
\textbf{El primer Gap} (mundo real vs. mundo observable): cada persona solo puede observar una parte ínfima del mundo, así que el conocimiento de cada persona tiene limitaciones enormes. La esencia de este gap es \textbf{lo limitado de la capacidad de observación}.

\textbf{El segundo Gap} (mundo observable vs. mundo simulado): «hay una distancia enorme». Pero la razón es distinta: este gap existe porque \textbf{no tenemos suficiente compute ni suficiente data} para modelar el mundo real. En algunos escenarios clásicos (como la correlación entre CFD y las pruebas de túnel de viento) ya es muy buena, pero «antes de que un avión despegue de todos modos hay que probarlo en el túnel de viento, no se puede decir que con el CFD ya basta».
\end{importantbox}

Su reflexión al respecto está cargada de sentido filosófico: «Es bastante desalentador. Tal vez nosotros, como especie, vivamos dentro del simulador de alguien más, y aun así fantaseamos con construir nuestro propio simulador capaz de simularnos a nosotros mismos». Pero al mismo tiempo considera que un modelo como GPT-4, capaz de mimic el funcionamiento del cerebro humano, ya «es un milagro».

\subsection{Video Model vs. Physical Simulation}

«¿Por qué seguimos usando Taichi para hacer rendering, para hacer physical simulation? ¿Por qué no usar Sora?» — esta pregunta apunta directo al valor de la existencia de la disciplina de Graphics en la era de la IA.

Hu Yuanming (胡渊鸣) respondió con mucho pragmatismo: «Hace cinco años, Frozen se escribía en 99 líneas; ahora, al prompt que se le da a Sora le basta con 20 palabras. En algunos escenarios de verdad se le acabó el oficio a la gente de Graphics. Pero muchas de las cosas que hace Graphics fueron las que hicieron posible que naciera Sora.»

La dirección final de esta evolución es «predominantemente learning», pero el simulator sobrevivirá de alguna forma: como inductive bias, como fuente de synthetic data, o como parte de un learnable module. Neural Graphics (DLSS, NeRF, etc.) ya está reemplazando gradualmente al ray tracing tradicional, y el posicionamiento de Taichi 2.0 es precisamente «supporting Graphics and AI».

\subsection{Sim-to-Real Gap: el reto fundamental de la simulación}

Hu Yuanming señala que todos los simulator existentes (Isaac, MuJoCo, Drake, PyBullet, Genesis) enfrentan un problema en común: el \textbf{sim-to-real gap}.

\begin{warningbox}{Las leyes físicas son exactas, pero las condiciones de frontera no}
La segunda ley de Newton es exacta, y la relatividad también lo es. Pero muchas boundary condition no lo son: por ejemplo, el modelo de fricción; el mecanismo de fricción de una superficie lisa y el de una superficie rugosa son completamente distintos, y la fricción de un vidrio sobre una superficie mojada es distinta otra vez. Todos estos modelados son inexactos.
\end{warningbox}

Él considera que los simulator del futuro necesitan volverse más data-driven y combinarse con redes neuronales. Esto no significa usar un Transformer para volver a aprender la segunda ley de Newton (no tiene sentido gastar miles de millones de FLOPS en algo que originalmente se puede hacer con unos cuantos FLOPS), sino dejar que los métodos data-driven compensen las partes que \textbf{los modelos aproximados no pueden describir con precisión}, como la stress-strain curve de los materiales.

\subsection{Inductive Bias: el valor del prior físico en la IA}

Hu Yuanming subrayó en particular la importancia del inductive bias (sesgo inductivo) en los sistemas de IA:

\begin{importantbox}{La eficiencia del 3D Inductive Bias}
Tomando la video generation como ejemplo: en la graficación tradicional, la camera projection solo necesita unos 50 FLOPS, pero si se le pide a la IA que aprenda ese mapeo, tal vez haya que agregar parámetros del orden de billion. Por eso, la camera projection debería quedar como hardcode, dejándole al modelo la capacidad para las partes que de verdad necesitan aprenderse. De forma similar, en un video model, cuando la camera gira 360 grados aparece un mundo completamente distinto; si existe una representación 3D básica, al menos se puede resolver el problema de physical persistency.
\end{importantbox}

Resumió esto en un principio de equilibrio muy incisivo:

\begin{quote}
«El simulator necesita tener más AI adentro, y la AI necesita tener más 3D inductive bias adentro. Hay que encontrar el punto de equilibrio exacto: no hay que usurpar funciones que no corresponden, no hay que subestimar la capacidad de aprendizaje del modelo, pero tampoco hay que dejarle todo al learning.»
\end{quote}

\subsection{Las limitaciones de Synthetic Data}

Sobre la resistencia de Sergey Levine hacia la simulation en su blog «The Spork of AGI» —donde sostiene que la simulation en realidad es una limitación para los robotics foundation model y que se debería usar real world data—, Hu Yuanming dijo que «tiene sentido»:

\begin{warningbox}{Synthetic Data no puede ser la fuente principal de datos}
«Si toda la data es synthetic data, o si la mayor parte de la data es synthetic data, definitivamente no va a work. Porque la data que se puede synthesize siempre tiene que follow ciertas rules, y esas rules necesariamente las pensó un ser humano. De cien millones de situaciones en el mundo real, las rules tal vez solo alcancen a cover diez mil.» — Pero synthetic data sí puede usarse para enhance un modelo ya existente, en lugar de build from scratch.
\end{warningbox}

\subsection{Los límites de la simulación y los universos anidados}

Al hablar sobre «los límites de la simulación física», Hu Yuanming dio una respuesta que invita a la reflexión. Si nuestro propio mundo ya es algo simulado, entonces \textbf{el simulador dentro del simulador no puede ser más potente que el simulador matriz}.

\begin{knowledgebox}{El argumento de complejidad computacional de la simulación anidada}
«Tal vez tú, dentro de nuestro mundo, construyas un simulador que simule el mundo entero y que se ejecute más rápido que el tiempo real; eso probablemente sea unlikely. A menos que solo hagas una simulación small scale. Esto también significa que \textbf{predecir el mercado de valores probablemente sea muy difícil}, porque el mercado de valores involucra a todas las personas del mundo entero participando en él: el sistema que necesitarías simular sería tan grande como el mundo real mismo. Pero hacer una simulación de un túnel de viento sí es posible.»
\end{knowledgebox}

Esta perspectiva conecta con elegancia la reflexión filosófica de la infancia sobre «si el árbol de durazno crece cuando no lo estoy viendo» con la teoría de la complejidad computacional de la vida adulta.

\subsection{La caverna de Platón de los LLM}

Cuando le preguntaron si «los LLM como GPT viven en la caverna de la que hablaba Platón», la respuesta de Hu Yuanming fue inesperada:

\begin{quote}
«Sí. Porque su loss function es su training data, y todo lo que hace es minimize esa loss function. Pero, ¿acaso cada uno de nosotros no vive también dentro de la caverna de Platón? Es posible que cada uno de nosotros vea incluso menos que un LLM. Solo que el LLM a veces comete errores de sentido común, y por eso uno se ríe de él seguido. Pero, ¿de verdad somos mejores que el LLM? Yo creo que no necesariamente.»
\end{quote}

Él considera que la AGI puede surgir a partir de la text, video e image data ya existente, y que no necesariamente hace falta interiorizar en ella todo el conocimiento humano sobre el mundo físico. Lo esencial es darle a la AGI la capacidad de \textbf{tool use}: dejar que tenga su propia computadora y que pueda escribir código para analizar fenómenos físicos. «Puede descubrir por sí misma su propia ley de Newton.»

\subsection{Si se hiciera de nuevo un Simulator}

Si hoy se volviera a hacer un simulator, Hu Yuanming lo impulsaría \textbf{predominantemente con learning}, y no con la ley de Newton como base. Los «elementos subyacentes» que conservaría incluyen:
\begin{itemize}
\item Los objetos no desaparecen de la nada (physical persistency)
\item Dos objetos no pueden ocupar la misma posición (collision)
\item La perspective projection de la camera
\end{itemize}

Y convertiría todas las physical law (incluidas la deformación elastoplástica, la dinámica de fluidos, etc.) en algo data-driven. El sistema resultante no sería ni un simulator tradicional ni una neural network pura, sino \textbf{«un 3D inductive bias con una neural network incrustada»}: un prior estructurado que le facilita el aprendizaje a la neural network.

\subsection{¿La AGI necesita física?}

Hu Yuanming considera que la AGI «tal vez no necesite en particular la ley de Newton ni las ecuaciones de Lagrange», pero sí necesita \textbf{tool use}: la AGI necesita tener su propia computadora, capaz de escribir código finite element para analizar fenómenos físicos.

\begin{knowledgebox}{La limitación esencial del Next Token Prediction}
Sin importar cuánta o cuán poca información contenga un token, cada token de un Transformer consume una cantidad fija de compute ($2 \times$ parameter count / sparsity via MoE). Pero, por naturaleza, algunos token necesitan más compute que otros; ese compute adicional no puede ocurrir dentro del paradigma actual del LLM, sino que requiere que el LLM opere herramientas externas. Por eso la AGI no necesita que se le interiorice todo el conocimiento humano sobre el mundo físico: puede «descubrir por sí misma su propia ley de Newton».
\end{knowledgebox}

Esta idea tiene una consecuencia importante: la forma actual de scaling up un Transformer (aumentar la cantidad de parámetros, para que cada token reciba más compute) es, en esencia, resolver por fuerza bruta un problema que existe a nivel de la arquitectura. Un enfoque más inteligente tal vez sea —como hace Taichi— procesar de manera distinta cómputos de granularidades distintas.

\subsection{Reflexiones sobre Robotics}

Aunque la formación técnica de Hu Yuanming está muy relacionada con Robotics (simulación física, programación diferenciable, generación 3D), él admite que su «nivel de interés en Robotics no es tan alto».

\begin{quote}
«Creo que Robotics es algo con un impact enorme. Tal vez la única razón por la que no me he puesto a hacer Robotics es que, hoy por hoy, mi nivel de interés en esto no es tan alto. Creo que solo puedo hacer bien las cosas que de verdad me gustan.»
\end{quote}

Pero también reconoce que existe un punto de intersección potencial entre Robotics y Meshy: muchos socios ya están usando Meshy para generar activos 3D como tazas, tazones y mesas para el entrenamiento de Robotics. Si en el futuro Meshy soporta articulated objects (objetos con articulaciones), esta intersección se volverá aún más notable.

Su juicio sobre el papel de la simulation en Robotics también es muy perspicaz: hacer que un robot manipule las verduras y los platos sobre una mesa \textbf{es más difícil que hacer volar un avión}: «porque allá arriba no hay más que aire, pero cuando necesitas manipulate algo sobre la mesa con más detalle, la inconsistencia entre las distintas condiciones de frontera hace que sea muy difícil llevar la simulation a la realidad.»

\subsection{Resumen de la sección}

La simulación física atraviesa un momento clave de transición de rule-based a data-driven. El cuello de botella de la simulación tradicional no está en las leyes físicas en sí, sino en el modelado preciso de las condiciones de frontera. La dirección futura es conservar el prior 3D subyacente (como inductive bias) y dejar que el learning se encargue de las partes complejas y difíciles de modelar con precisión. Taichi 2.0 está evolucionando en esa dirección.
\section{Meshy AI: la práctica de comercialización de la generación 3D}

\subsection{Del PhD a la decisión de emprender}

Al graduarse, Hu Yuanming (胡渊鸣) enfrentaba varias opciones: (1) ir a una escuela para ser profesor; (2) trabajar en una gran empresa; (3) emprender.

Sobre trabajar en una gran empresa, su respuesta fue tajante: «Definitivamente no lo soportaría. No porque una gran empresa sea cómoda o mala, sino porque no soporto que alguien me diga qué hacer.» Durante su paso por el MIT, su asesor solo se reunía con él media hora a la semana (menos aún durante la COVID), pero la ayuda que le daba era enorme: aliento, orientación de dirección, presentaciones con gente destacada. «Pero en lo técnico específico, él no me decía qué hacer.»

\begin{knowledgebox}{La psicología del «no lo soporto» del fundador técnico}
La motivación de Hu Yuanming para emprender no era el dinero («posiblemente solo represente entre el 5\%--10\%»), sino una necesidad extrema de autonomía. Este rasgo psicológico se observa en muchos fundadores técnicos exitosos: no pueden trabajar dentro de un marco ya establecido y tienen que definir por sí mismos el problema y la dirección. Es un arma de dos filos: impulsa a hacer lo que otros no se atreven a hacer, pero también implica que, al inicio de una empresa, faltarán muchas de las estructuras y la experiencia que tendría una «empresa normal».
\end{knowledgebox}

Sobre la vía académica: «El principal problema es que la perspectiva académica no era suficientemente amplia: no me había dado cuenta de que existía algo tan bueno como GenAI para hacer.» Hacia el final del doctorado, el campo de Graphics le parecía haber llegado a un límite, y no sabía qué más podría hacer en el mundo académico. «Así que pensé: ¿podría convertir en un producto comercial lo que había hecho en el doctorado?»

\subsection{La transición de Taichi Graphics (太极图形) a Meshy}

Tras doctorarse en 2,021, Hu Yuanming fundó la empresa Taichi Graphics, y durante los primeros 18--24 meses estuvo explorando la dirección de comercialización. Taichi agregó soporte para más hardware (GPU de AMD, backend de JavaScript) y probó algunas formas de comercialización, pero al final descubrió que \textbf{solo podía hacer subcontratación}: desarrollar un simulator para otros y entregarlo, con un margen bruto muy bajo, y el mercado 2B chino resultaba aún más insostenible.

\begin{warningbox}{La trampa de comercialización del emprendimiento tecnológico}
«En esa época había varias palabras muy de moda —metaverso, código abierto, infrastructure— y uno no podía evitar pensar "¿será que lo que estoy haciendo puede servir para algo dentro de esto?" Es algo muy difícil de evitar, pero cuando ya tienes más experiencia comercial, dejas de caer en esos errores básicos.» — dejarse arrastrar por las palabras de moda de la industria, en lugar de volver a los primeros principios y pensar en qué es lo que el usuario realmente está dispuesto a pagar.
\end{warningbox}

En este proceso llegó a una idea clave sobre su propia personalidad: «Un Founder con una personalidad como la mía solo puede hacer productos estandarizados orientados al consumer, porque odio tratar con grandes clientes.» Esto determinó que, desde el inicio, el modelo de negocio de Meshy fuera subscription + API, en lugar de personalización a nivel empresarial.

El punto de inflexión clave llegó con el pensamiento de «primeros principios» del segundo semestre de 2,022:

\begin{importantbox}{La lógica comercial de Meshy}
\begin{itemize}
\item ChatGPT puede generar lenguaje
\item Stable Diffusion puede generar imágenes
\item La gente va a sitios como Sketchfab y \textbf{paga por comprar} modelos 3D
\item Pero en el mercado \textbf{todavía no había} un producto capaz de generar modelos 3D
\end{itemize}
Conclusión: «Si podemos generar modelos 3D a un costo cercano a cero — then we are rich.»
\end{importantbox}

Detrás de esta transición está la clásica historia de Intel al pasar de la DRAM a la CPU: ante la presión competitiva de los semiconductores japoneses, Andy Grove y Gordon Moore propusieron «¿por qué no nos despedimos a nosotros mismos, salimos de esta sala y volvemos a entrar, y pensamos qué haría un nuevo CEO?», y al final decidieron pasar de la DRAM a la CPU. Tres años de una transición dolorosa dieron lugar a varias décadas de éxito posterior de Intel.

Hu Yuanming reconoce abiertamente que la transición encontró resistencia: «También hubo personas que, al no interesarles el tema, dejaron la empresa. Pero quienes se quedaron y siguieron adelante ahora consideran que fue una decisión muy acertada.» De hecho, Meshy es el tercer producto de la empresa: el primero fue el propio Taichi (que fracasó comercialmente), el segundo fue otro intento (que también se canceló), y al final fueron all in con Meshy.

El punto de partida técnico de Meshy vino del ecosistema de Taichi: alguien había implementado Gaussian Splatting con Taichi, alguien más había implementado NeRF, y el equipo pensó: «Resulta que Taichi también puede hacer esto, ¿no podríamos empujar un poco más en la dirección de generar activos 3D?»

\begin{importantbox}{La lógica central detrás del Pivot de una empresa emergente}
«Con franqueza, en ese momento Meshy era la mejor dirección de Pivot. Porque si me hubiera puesto a hacer LLM, 2D Diffusion, Text-to-Image, Text-to-Video, no necesariamente le habría ganado a otros. Las limitaciones de mi propia capacidad hicieron que Meshy fuera, en ese momento, mi mejor dirección.» — reconocer los límites del círculo de competencia propio y del equipo, y elegir la mayor oportunidad de mercado dentro de ese círculo.
\end{importantbox}

\begin{importantbox}{Despedirme a mí mismo cada tres o seis meses}
«En el proceso de dirigir la empresa, cada tres a seis meses tengo que despedirme a mí mismo. Si uno quiere sobrevivir en un entorno de emprendimiento tan intense, lo mejor es pensar que el uno mismo de hace tres meses era un idiota; eso es lo mejor. Si no piensas que el tú de hace tres meses era un idiota, significa que en esos tres meses no creciste.»
\end{importantbox}

\subsection{El lanzamiento en 8 horas y el juego de terror de las cuatro caras}

La primera versión de Meshy se lanzó en tan solo 8 horas y de inmediato consiguió 1,000 usuarios. Pero la calidad de generación en ese momento era pésima: el método de texturizado consistía en generar una textura para cada uno de los lados —frente, atrás, izquierda y derecha— del modelo y pegarlas, lo que hacía que la cabeza del personaje tuviera 4 caras.

«Cuando hacía entrevistas con los usuarios, siempre les decía: por ningún motivo le digan a nadie que Meshy lo hizo Hu Yuanming.» Solo después de las iteraciones de Meshy-1 y Meshy-2 se atrevió a reconocer públicamente el producto.

\begin{knowledgebox}{De las cuatro caras al prototipo de videojuegos 3A}
Al principio, Meshy solo servía para hacer juegos de terror (por las 4 caras), pero con la iteración constante de la tecnología, ahora hay estudios de videojuegos AA e incluso AAA que usan Meshy para hacer prototipos. «No hay que subestimar la velocidad con la que avanza la tecnología» es la mayor lección que Hu Yuanming aprendió en este proceso. En 2,023 dijo que se había resuelto el 10\%; visto en retrospectiva desde 2,025, se resolvió cerca del 90\% de lo que se imaginaba en ese momento, pero el límite superior pasó del 100\% al 200\%, así que en realidad es alrededor del 50\%.
\end{knowledgebox}

\subsection{Modelo de negocio y posicionamiento de mercado}

El modelo de negocio de Meshy es, en esencia, el mismo que el de OpenAI: \textbf{suscripción + API}. Los datos centrales actuales:

\begin{table}[H]
\centering
\caption{Datos comerciales centrales de Meshy}
\begin{tabular}{ll}
\toprule
\textbf{Indicador} & \textbf{Dato} \\
\midrule
Usuarios registrados & Más de 4 millones \\
Visitas mensuales al sitio & Cerca de 2.5 millones \\
Cliente más grande & Meta (según el valor promedio por cliente) \\
Participación de mercado en Estados Unidos & Cerca del 55\% \\
Comparación con la competencia & Aproximadamente igual a la suma de las visitas del segundo y el tercer lugar \\
Crecimiento de ingresos mensual & Cerca del 20\% \\
Crecimiento de ingresos anual & Más de 10 veces \\
\bottomrule
\end{tabular}
\end{table}

En cuanto al posicionamiento de mercado, la generación 3D es un mercado de \textbf{tamaño mediano}: más pequeño que el de los video model (lo que evita competir de frente con las grandes empresas), pero más grande que muchos nichos verticales (lo suficiente para sostener una empresa emergente). El potencial del mercado de consumo es igualmente enorme: en Estados Unidos, los envíos anuales de impresoras 3D crecen 20\%, y muchas personas que compraron una impresora no tienen modelos que imprimir.

\begin{warningbox}{La selección de mercado de una empresa emergente}
«Si el mercado en el que se trabaja es demasiado grande (como el de los video model), el reto es que no puedes ganarles en la competencia feroz a las grandes empresas; si el mercado es muy pequeño, no alcanza para sostener a la empresa. Encontrar un mercado de tamaño mediano permite evitar competir con los gigantes y, al mismo tiempo, crecer muy rápido.»
\end{warningbox}

La visión de más largo plazo de Hu Yuanming para Meshy va más allá de los modelos 3D: «En esencia, lo que estamos haciendo es usar multimodal AI para traerle diversión a la gente: AI for fun.» Si se logra resolver este planteamiento más amplio, el espacio de mercado de Meshy superará con mucho el ámbito actual de lo 3D. Su meta es construir una empresa como NVIDIA.

\subsection{Por qué fue un equipo de Graphics el que logró esto}

Un hecho contraintuitivo es que, antes de esto, casi nadie en el equipo de Meshy había entrenado un modelo grande: eran un grupo de personas con formación en Graphics.

\begin{knowledgebox}{La ventaja interdisciplinaria de la gente de Graphics}
«La gente que se dedica a Graphics tiene una característica muy marcada: son sumamente sólidos técnicamente, aprenden lo que sea. Si todos los días trabajas con las ecuaciones de Euler-Lagrange, que hoy te pongan a aprender Diffusion es facilísimo. Si te la pasas trabajando con MCMC, resolver hoy una SDE es algo que está totalmente relacionado.» Esta sólida base matemática y de ingeniería hace que las personas con formación en Graphics se muevan como pez en el agua en la era de la AI.
\end{knowledgebox}

\subsection{La evolución técnica de Meshy: de NeRF a 3D Native Generation}

La ruta técnica de Meshy pasó por una evolución notable. La tecnología usada al principio era «muy rudimentaria»: primero se usaba 2D diffusion para generar, a partir de texto, cuatro imágenes (frente, atrás, izquierda y derecha), que luego se pegaban sobre el modelo 3D (de ahí viene el «juego de terror de las cuatro caras»). Pero con la iteración de la tecnología, ahora ya se ha llegado a 3D native generation.

\begin{table}[H]
\centering
\caption{Evolución técnica y de producto de Meshy}
\begin{tabular}{llll}
\toprule
\textbf{Versión} & \textbf{Tecnología central} & \textbf{Nivel de calidad} & \textbf{Caso de uso} \\
\midrule
Etapa inicial & Texturizado 2D multivista & Pésimo (cuatro caras) & Juegos de terror \\
Meshy 1 & NeRF + texturizado mejorado & Utilizable & Prototipos de juegos independientes \\
Meshy 2 & 3D Native + Diffusion & Bastante bueno & Prototipos de juegos AA \\
Actual & Pipeline de AI multimodal & Bueno & Algunos juegos AAA \\
Dirección futura & Articulated objects & --- & Entrenamiento de Robotics \\
\bottomrule
\end{tabular}
\end{table}

Reconoce abiertamente que Meshy todavía no llega al nivel de dejar «totalmente satisfechos» a los estudios de videojuegos: «Con franqueza, ahora mismo la producción de modelos 3D con AI todavía no llega de verdad al punto en que un studio piense "¡wow, estoy muy satisfecho!"; casi siempre todavía tienen que retocarlo con ZBrush.» Pero el ritmo de avance es sumamente rápido: «No hay que subestimar la velocidad con la que avanza la tecnología» es la lección que confirmó una y otra vez en este proceso.

Compartió una idea clave sobre los usuarios: cuando antes fracasó la comercialización de Taichi, tras hablar con muchos usuarios descubrió: «Todos decían "no voy a pagar por este software tuyo, pero si me vendes todos los modelos que hay dentro de él, sí pagaría por ellos".» Esta idea dio origen directamente al modelo de negocio de Meshy.

\subsection{La estrategia competitiva de Meshy: una reflexión comparativa con Cursor}

Cuando le preguntaron si era «un buen CEO», Hu Yuanming ofreció un análisis comparativo franco:

\begin{knowledgebox}{Buen sector vs. buena ejecución}
«¿Dices que Meshy crece rápido? Comparado con Cursor y Lovable, ni siquiera se puede llamar rápido. Con un solo producto de AI coding, ellos llegan a 200M ARR. Nosotros somos muy hardworking, ellos también son muy hardworking, pero el sector de ellos tiene un retorno más alto.» Esto llevó a una reflexión profunda: para un CEO, \textbf{elegir el sector} puede ser más importante que la ejecución. «Pero lo que se obtiene por suerte, tarde o temprano hay que devolverlo con méritos propios; que Cursor sea rápido ahora también es algo pasajero, lo importante es ver la siguiente etapa.»
\end{knowledgebox}

También citó una comparación aleccionadora: Unitree Robotics (宇树科技), el fabricante chino líder de robots humanoides, tiene ingresos de más de mil millones de yuanes, pero eso es solo una décima parte de los de Pop Mart (泡泡马特), la marca de juguetes de moda. «Esto muestra que la diferencia de retorno entre sectores distintos es enorme: Andrej Karpathy tiene razón, algo que rinde diez veces más muchas veces solo es dos veces más difícil.»

Pero también ve el asunto de manera dialéctica: «Estás en la mesa de juego, tienes un buen equipo, tienes un negocio que crece de manera constante, tienes la oportunidad de hacer cosas nuevas: siempre existe la posibilidad de construir una empresa como NVIDIA, siempre que tengas ese gen.»

\subsection{Oportunidades y desafíos de la era de la AI}

Cuando le preguntaron qué oportunidades y desafíos implica hacer Meshy en la era de la AI, Hu Yuanming respondió desde dos planos: el técnico y el comercial:

\textbf{Oportunidades}:
\begin{itemize}
\item Las industrias upstream y downstream han ido madurando: Text-to-Image (Black Forest Lab), Language Generation (ChatGPT, Claude), Video (Gemini/Veo); muchas cosas ya no requieren que uno mismo haga build, basta con conectar una API
\item Tecnologías como Diffusion, RF, DiT y MoE han ido madurando, lo que hace posible construir modelos de generación 3D
\end{itemize}

\textbf{Desafíos}:
\begin{itemize}
\item El ritmo de la época es sumamente rápido; hay que actualizar el producto con mucho esfuerzo para mantener una buena posición
\item La empresa puede verse desplazada por nuevas tecnologías, pero para una compañía con una gran capacidad de adaptación, «el espacio de mercado es infinito, porque siempre se pueden resolver nuevos problemas»
\end{itemize}

Sobre la intensidad de su trabajo dice que «no es muy apropiado contarla en el programa; si no, a lo mejor la gente ya no se atrevería a emprender». Pero cada día, al levantarse, siente que está «good enough, y todavía puede recibir los nuevos retos con toda la energía». Cuando le preguntaron si se cansaba más durante el PhD o ahora, no dudó: «Claro que ahora me canso más. Durante el PhD la pasaba increíble: solo tenía que cuidarme a mí mismo. Pero ahora toda la empresa tiene que ganar, y todos tienen que poder comer.»

\subsection{Las dificultades de la industria de los videojuegos y la oportunidad de Meshy}

Tras tratar con muchas empresas de videojuegos, Hu Yuanming tiene un conocimiento profundo del estado actual de esa industria. En los últimos años, toda la industria se encuentra en contracción (un mercado de alrededor de 300,000 millones de dólares) y enfrenta múltiples desafíos:

\begin{enumerate}
\item \textbf{Competencia por la atención}: plataformas de video corto como TikTok le quitan a los usuarios el tiempo que antes dedicaban a jugar
\item \textbf{Estancamiento en la jugabilidad}: desde PUBG (comer pollo), durante mucho tiempo no ha surgido una nueva jugabilidad realmente disruptiva
\item \textbf{Espiral de costos de producción}: GTA necesita un ciclo de producción de 10 años, los costos son cada vez más altos, pero el precio de los juegos está limitado (un juego de 20 horas solo se puede vender en 30 o 40 dólares, mientras que con ese mismo precio una película solo dura entre 2 y 3 horas)
\item \textbf{Alta proporción de costos en activos 3D}: el 50\% del costo de producción de un juego 3A está en los 3D Asset
\end{enumerate}

«Por eso tienen muchas ganas de reducir costos y aumentar la eficiencia, y necesitan algo como Meshy.»

Sobre el mercado de consumo de la generación 3D, dio un ejemplo muy gráfico: en su familia, tres generaciones (sus padres, él mismo y su novia) usan LLM, pero no ha podido convencer a sus padres de usar Meshy: «El 3D es un mercado relativamente de nicho; si mandas un archivo OBJ por WeChat, no se puede abrir.» Pero si se logra hacer una reconstrucción 3D de personas de alta calidad (por ejemplo, convertir una foto de grupo en un modelo imprimible en 3D), «todo el mundo sería usuario de Meshy».

\begin{knowledgebox}{La lógica de selección de mercado de un Startup}
«Si el mercado es demasiado grande (como el de los video model), no puedes ganarle a las grandes empresas en la competencia feroz; si el mercado es demasiado pequeño, no alcanza para sostener a la empresa. Hay que buscar un mercado de tamaño mediano, uno que a las grandes empresas no les interese hacer, pero que sea suficiente para que una empresa emergente eche raíces y se desarrolle.» Meshy está justo en ese punto óptimo: el mercado de la AI generativa 3D es lo bastante grande para permitir que la empresa alcance una escala de salida a bolsa, pero no tanto como para atraer la competencia total de todas las grandes empresas.
\end{knowledgebox}

\subsection{De Taichi a Meshy: el salto interdisciplinario del equipo de Graphics}

Un hecho contraintuitivo es que, antes de esto, casi nadie en el equipo de Meshy había entrenado un modelo grande: eran un grupo de personas con formación en Graphics, y la pila tecnológica de Meshy casi no tiene relación directa con la de Taichi. Esto equivale a empezar desde cero en un campo completamente nuevo.

\begin{knowledgebox}{La ventaja interdisciplinaria de la gente de Graphics}
«La gente que se dedica a Graphics tiene una característica muy marcada: son sumamente sólidos técnicamente, aprenden lo que sea. Si todos los días trabajas con las ecuaciones de Euler-Lagrange, que hoy te pongan a aprender Diffusion es facilísimo. Si te la pasas trabajando con MCMC, resolver hoy una SDE es algo que está totalmente relacionado.» Esta sólida base matemática y de ingeniería hace que las personas con formación en Graphics se muevan como pez en el agua en la era de la AI.
\end{knowledgebox}

Cuando le preguntaron si no le preocupaba que el equipo no fuera adecuado para hacer text-to-3D, su respuesta fue muy franca: «No tengo ninguna preocupación al respecto, porque en ese momento nadie estaba haciendo esto.» Esta es la first-mover advantage de Meshy: entrar cuando no hay consenso, y convertir la curva de aprendizaje en una barrera competitiva.

\subsection{Resumen de la sección}

El éxito de Meshy confirma varias decisiones clave: (1) entrar cuando no hay consenso: en ese momento nadie hacía generación 3D, y más tarde se convirtió en el primer producto público en lanzarse; (2) la capacidad de adaptación de un equipo con genes técnicos: un equipo con formación en Graphics que aprendió rápido la tecnología de AI; (3) intuición comercial: los usuarios están dispuestos a pagar por modelos 3D, así que se satisface esa necesidad generándolos a un costo cercano a cero. Del «si algo fracasa, es como si no se hubiera hecho» al primer lugar del mercado, lo central fue una velocidad de iteración sumamente rápida y una comprensión precisa de las necesidades de los usuarios.
\section{La autoevolución del CEO: valentía y sabiduría}

\subsection{Visión del talento: Hungry, Humble, Smart, Clarity}

Las cuatro características que Hu Yuanming (胡渊鸣) más valora al entrevistar a cada candidato de Meshy:

\begin{importantbox}{El modelo de cuatro dimensiones del talento de Hu Yuanming}
\begin{enumerate}
\item \textbf{Hungry}: aún no se ha demostrado a sí mismo, tiene muchas ganas de convertirse en una mejor versión de sí mismo, siente una especie de "hambre" por generar impact. Quien ya alcanzó el éxito y el reconocimiento puede no estar dispuesto a take risk.
\item \textbf{Humble}: es capaz de reconocer que todavía tiene muchas cosas que no ha hecho bien, tiene growth mindset, y se transforma y se cuestiona a sí mismo constantemente. "Si una persona no es humble, por muy bueno que sea en este momento, en el futuro no avanzará."
\item \textbf{Smart}: el mejor problem solver, aquel que, con recursos limitados, encuentra todas las formas posibles de resolver un problema (por ejemplo, entrenar con números de punto flotante de menor precisión, o encontrar un mejor optimizer).
\item \textbf{Clarity}: expresarse con claridad. "Cuando te das cuenta de que no entiendes en absoluto lo que dice una persona, en nueve de cada diez casos el problema no eres tú, sino que esa persona no ha pensado las cosas con claridad." Expresarse de forma poco clara casi siempre no es un problema de elocuencia, sino de pensamiento poco claro.
\end{enumerate}
\end{importantbox}

Él admite que también comete errores frecuentemente en estos cuatro aspectos. Uno de los cambios más grandes tiene que ver con Clarity: pasar de "decir mucho pero que solo el 10\% sea útil" a darse cuenta de que "escuchar es mucho más importante que hablar". En cuanto a Smart, el error más grande fue "be smart about the solution, but ignore being smart about strategy": pasar demasiado tiempo pensando en el how y muy poco tiempo pensando en el why.

\subsection{La brecha cognitiva entre ser Intern y ser CEO}

Toda la experiencia profesional de Hu Yuanming antes de ser CEO consistió en haber sido becario en NVIDIA, Adobe y Microsoft. "Mi experiencia más professional fue haber sido intern: eso era todo mi conocimiento del mundo laboral antes de ser CEO."

El struggle inicial no era "cómo ser un buen CEO", sino "ni siquiera sé cómo debería ser una empresa normal". Además, como el financiamiento de Taichi (太极) fue muy fluido, no existía la presión de sobrevivencia de "el próximo mes no podré pagar los salarios", y esto en realidad delay el momento en que tuvo que enfrentar los hechos crueles: "que el financiamiento fluya bien en realidad es un arma de doble filo; te hace sentir menos dolor, y entonces no sientes suficiente urgent para cambiar".

\begin{knowledgebox}{Un PhD no es condición suficiente para ser CEO}
"Un PhD te puede enseñar a ser humble, work hard, presentation skills, tecnología de vanguardia. Pero el mercado, los usuarios, la producción en volumen de un producto: eso el PhD no lo toca." Descubrió que ejecutivos de grandes empresas veinte años mayores que él, que salieron a emprender, "parece que lo hicieron incluso peor que nosotros"; esto demuestra que no es un problema de experiencia, sino de velocidad de aprendizaje y de qué tan humble se es. Incluso Li Xiang (李想), de Li Auto (理想汽车), que ni siquiera terminó la preparatoria, lo hizo muy bien: "un PhD no es condición necesaria, la visión sí lo es."
\end{knowledgebox}

Se autoevalúa como "so far un CEO que apenas aprueba, a duras penas", y entre los aspectos que necesita mejorar están: delegate mejor el trabajo, administrar el tiempo, ser más decidido en ciertas cosas, aprender más rápido, dedicar más tiempo a lo externo que a lo interno, leer más libros, reflexionar durante más tiempo, y pensar en la visión de la empresa.

\subsection{La cultura de ganar y la cultura de perder}

\begin{warningbox}{Características de la cultura de perder}
\begin{itemize}
\item \textbf{Desgaste interno}: los departamentos se desgastan entre sí en lugar de cooperar
\item \textbf{Pasar la responsabilidad a otros}: cuando surge un problema, no se revisa lo ocurrido ni se reflexiona
\item \textbf{Poco realista}: no se trata de que la meta sea alta (la meta puede seguir siendo alta), sino de no enfrentar los hechos: disfrazar lo malo como si fuera bueno
\item \textbf{Buscar la armonía}: el equipo no puede criticar al CEO, no puede señalar los problemas en público
\end{itemize}
Todo esto termina llevando a la empresa a "perder".
\end{warningbox}

Por el contrario, el pilar cultural de Meshy es "Respect Facts" (respetar los hechos). Hu Yuanming dio un ejemplo concreto: hace unos días, dos colegas lo "llevaron al pasillo" para criticarle problemas de gestión. "Cuando lo escuché se me encogió el corazón, pero dejé que terminara de hablar, y luego pensé si de verdad no lo había hecho bien; después descubrí que efectivamente había cosas que no había hecho bien."

Él considera que poder contratar a este tipo de personas y tener este tipo de ambiente ya es en sí una suerte: "Esa persona seguramente tiene un fuerte sentido de ownership hacia la empresa, se ve a sí misma casi como el CEO, y por eso está dispuesta a decirte estas cosas."

En una empresa impulsada por la innovación, el CEO no necesariamente sabe todo: esta es la diferencia esencial con una empresa manufacturera. "En una empresa impulsada por la innovación tienes que aceptar que el CEO no necesariamente sabe todo."

\subsection{La experiencia de trabajar en Meshy}

Empleados veteranos con tres o cuatro años de antigüedad describen a Meshy como "la mejor empresa en la que he trabajado en más de diez años de carrera". Las razones incluyen: poder crecer rápidamente, tener un leader que escucha sus necesidades, desempeñar un papel importante en un negocio de rápido crecimiento, y tener acceso a la tecnología más de vanguardia (antes, en una gran empresa, tal vez solo escribían shaders, mientras que en Meshy pueden entrenar modelos con 512 GPU).

"En este proceso, todos se han convertido en una mejor versión de sí mismos." Esto coincide en gran medida con su objetivo como CEO: hacer que Meshy sea "un lugar que la gente más talentosa considere un excelente lugar para trabajar".

\subsection{El cambio de perspectiva más profundo: dejar de buscar que todos lo quieran}

Hu Yuanming dice que este fue su mayor momento de colapso cognitivo como CEO:

\begin{warningbox}{Buscar que todos te quieran es un acto egoísta}
"Si el objetivo de tu trabajo es que todos te quieran a ti, al final terminarás siendo odiado por todos." Porque el CEO tiene que tomar tough decisions: cancelar proyectos que le gustan a todos, despedir a colegas sin contribución pero con buenas relaciones, elegir entre el grupo de usuarios A y el grupo B (si haces lo que le gusta a A, B te critica). "Buscar que todos te quieran es poner lo personal antes que lo colectivo: cuando ocupas la posición de CEO, tienes que responder por los intereses de toda la empresa, no por que todos te quieran a ti."
\end{warningbox}

Citó la evaluación que Li Xiang, fundador de Li Auto (理想汽车), hizo sobre Elon Musk: "Evaluar a un CEO no es ver qué tan diplomático es, sino ver si en un momento de verdadera crisis es capaz de ver la esencia del problema de un vistazo."

Este cambio también se extiende al nivel de gestión. Descubrió que los nuevos managers suelen tener el mismo problema: como fueron tratados con consideración cuando ellos mismos eran IC (individual contributor), tratan a sus subordinados de la misma manera y no se atreven a exigir estándares más altos. La solución es contarles la "historia de sangre y lágrimas": "El problema que tiene hoy la empresa es porque hace medio año fui indeciso y quise que todos me quisieran."

\subsection{Las reglas de supervivencia de una startup}

\textbf{¿Qué es un buen negocio?} Hu Yuanming resumió tres dimensiones:
\begin{enumerate}
\item \textbf{Mercado grande}: la cuestión de si vale la pena hacerlo
\item \textbf{No consenso}: la cuestión de si todavía queda una oportunidad para ti (el timing es muy importante)
\item \textbf{Ventaja competitiva}: la cuestión de si tú mismo puedes lograrlo
\end{enumerate}

\begin{importantbox}{El volante tecnológico (Flywheel)}
Buena tecnología $\rightarrow$ buen producto $\rightarrow$ buen revenue $\rightarrow$ atrae a buena gente $\rightarrow$ produce buena tecnología. Si la tecnología no se puede comercializar, se convierte en "un espectáculo de fuegos artificiales que desaparece en cuanto termina". El volante de Meshy ya está girando: los ingresos mensuales crecen 20\%, cuenta con 4 millones de usuarios, y la tecnología se sigue iterando constantemente.
\end{importantbox}

\textbf{Sobre elegir una dirección}, citó a Andrej Karpathy: "Algo que rinde diez veces más suele ser solo el doble de difícil." Elegir la dirección correcta es mucho más importante que el nivel de esfuerzo: el equipo de Cursor es tan hardworking como el de Meshy, pero el sector en el que compite Cursor tiene un retorno mayor. Sin embargo, él también lo ve de manera dialéctica: "Lo que se obtiene por suerte, tarde o temprano hay que devolverlo con capacidad real. Que Cursor crezca rápido ahora es algo temporal; que en la siguiente etapa pueda crecer aún más dependerá de si el CEO es capaz de seguir aceptando nuevos retos."

\subsection{Respect Facts: el pilar cultural de respetar los hechos}

\begin{warningbox}{El hecho más cruel al emprender}
"La gran mayoría del esfuerzo en el mundo no tiene ninguna recompensa. Mientras más grande es lo que haces y mayor es el riesgo, muchas veces pones 99 puntos de esfuerzo y no obtienes ni un punto de recompensa. Pero tienes que hacer ese esfuerzo, porque no sabes cuál de esos 100 puntos es el que sirve. 'El cielo recompensa a los diligentes', 'un esfuerzo, una cosecha': todo esto es pura mentira."
\end{warningbox}

En la práctica de evaluación de Meshy, exigen pruebas ciegas: muestran de forma aleatoria la versión nueva y la antigua, sin decir cuál es cuál, y un tercero las evalúa. "Muchas veces, después de mucho esfuerzo, descubres que la versión anterior era mejor." Si no se enfrenta este hecho de manera objetiva, el equipo cae en el autoengaño.

\subsection{Bias for Action y empezar con el fin en mente}

Hu Yuanming se describe a sí mismo como alguien "muy bias for action": si quiere hacer algo, lo hace de inmediato, incluso sin dormir por la noche. Con frecuencia, cuando un proyecto todavía es apenas una idea sin forma, ya está pensando en el plano del futuro: "es un poco como Elon Musk construyendo cohetes, que antes de que el motor esté listo ya está pensando qué va a comer en Marte."

Pero esto no es un activismo ciego: él piensa "empezando con el fin en mente", es decir, qué debe hacer ahora para lograr el objetivo final, y luego planea hacia atrás a partir de ahí. Tanto Taichi como Meshy se planearon de esta manera, hacia atrás.

\begin{knowledgebox}{¿Si Jensen Huang volviera a empezar, no fundaría NVIDIA?}
"Creo que eso es alardear fingiendo modestia. Tienes una empresa de 4 billones de dólares; puedes decir lo que quieras, todo lo que digas será correcto, tú eres quien manda." Hu Yuanming cambió de tono: "\textbf{Yo sí volvería a elegirlo.} Aunque cada día es muy suffer, en este proceso se puede aprender mucho y convertirse en una mejor versión de uno mismo. Se ve una visión mucho más amplia que la del pequeño comedor del mundo académico."
\end{knowledgebox}

"En tiempos de paz, no hay nada que forje más a una persona que emprender." Esta es su evaluación final sobre el emprendimiento. No se puede esperar aprender demasiado en un entorno protegido (la escuela, una gran empresa), a menos que tengas un mentor excepcionalmente bueno. Emprender es la forma de crecimiento más cruel, pero también la más eficiente.

\subsection{Resumen de la sección}

La evolución personal del CEO es un proceso continuo de "matar a la versión pasada de uno mismo". Los cambios centrales incluyen: pasar de buscar que todos lo quieran a buscar el éxito del negocio; pasar de smart about solution a smart about strategy; pasar de trabajar con la cabeza agachada a abrir la mirada y comunicarse más. Las características de una cultura que puede ganar son enfrentar la realidad, aceptar la crítica y buscar la excelencia; las características de una cultura que puede perder son el desgaste interno, pasar la responsabilidad a otros y disfrazar los problemas.
\section{Valentía y sabiduría: Brain, Guts, Heart, Taste}

\subsection{Las tres dimensiones de Jensen Huang y la cuarta dimensión de Hu Yuanming}

El título del último capítulo de la entrevista, «Valentía y sabiduría», proviene del título de un artículo de Hu Yuanming en Zhihu. Ahí propone cuatro palabras clave para evaluar el talento: Brain, Guts, Heart, Taste; las tres primeras provienen de la respuesta de Jensen Huang en una conferencia en Stanford, y la cuarta la agregó él mismo.

\begin{importantbox}{{Brain / Guts / Heart / Taste}}
\begin{itemize}
\item \textbf{Brain} (mente): el problem solver más fuerte.
\item \textbf{Guts} (valentía): puede actuar rápido cuando la información es incompleta, puede aceptar el riesgo de fracasar, puede aprender del fracaso.
\item \textbf{Heart} (corazón): ser bondadoso desde el fondo; pero eso no significa ser la persona más nice. «A veces exigirle a tu equipo un estándar alto es, en realidad, la mayor bondad, porque si no marcas ese estándar el equipo va a fracasar.» En el fondo no puede haber intención de dañar a otros; debe ser una persona con principios y honesta.
\item \textbf{Taste} (gusto): la dimensión que Hu Yuanming agregó por su cuenta. Al contratar pregunta: «¿A qué juegos juegas? ¿Qué haces los fines de semana? ¿Qué libros lees?»; solo las personas con gustos e intereses afines pueden colaborar de cerca durante mucho tiempo.
\end{itemize}
\end{importantbox}

Sobre Taste, reveló que últimamente ha estado leyendo un libro de Lei Jun (雷军). Piensa que Lei Jun es un emprendedor del que vale mucho la pena aprender: en su momento fue un «trabajador modelo» en Zhongguancun, haciendo para Microsoft un sustituto de Office (la suite Pangu), y se perdió por completo la era de internet; más tarde resumió la teoría del cerdo volador: «en la boca del huracán, hasta los cerdos pueden volar». Esto sigue la misma lógica que su propia idea de que «elegir la dirección importa mucho más que esforzarse».

\subsection{Etapa uno y etapa dos: de quien resuelve problemas a quien los plantea}

Hu Yuanming propone un marco profundo para el desarrollo de la vida: \textbf{la etapa uno} y \textbf{la etapa dos}:

\begin{table}[H]
\centering
\caption{Diferencia esencial entre la etapa uno y la etapa dos}
\begin{tabular}{lll}
\toprule
\textbf{Dimensión} & \textbf{Etapa uno} & \textbf{Etapa dos} \\
\midrule
Origen del problema & Otros te plantean el problema & Tú mismo te planteas el problema \\
Claridad del objetivo & Hay un objetivo claro & No hay un objetivo claro \\
Alcance de la responsabilidad & Hacer bien lo que te corresponde & Pensar qué debe hacer todo el equipo \\
Habilidad central & Obedecer órdenes, competir al máximo por la eficiencia, perfeccionar habilidades concretas & Definir la dirección, enfrentar la incertidumbre \\
Criterio de éxito & Completar la tarea & Crear valor \\
\bottomrule
\end{tabular}
\end{table}

Este marco aplica a cada salto de la vida: de estudiante de licenciatura a estudiante de posgrado, de IC (colaborador individual) a Manager, de empleado de una empresa a emprendedor.

\begin{warningbox}{Cuanto más éxito tengas en la etapa uno, más puedes suffer en la etapa dos}
«Las habilidades de la etapa uno y las que exige la etapa dos son completamente distintas. En la etapa uno puede que se trate de obedecer órdenes, competir al máximo por la eficiencia, perfeccionar habilidades concretas. Pero al llegar a la etapa dos ya nadie te da órdenes, y ni siquiera se puede hablar de eficiencia, porque ni siquiera está definido el metric de la eficiencia. Mucha gente a la que le fue muy bien en la etapa uno se siente sumamente incómoda al llegar a la etapa dos: es como pelear contra un boss; terminas una barra de vida, el boss se transforma y empieza una segunda barra de vida, y si sigues peleando con los métodos de la primera etapa, mueres muy rápido.»
\end{warningbox}

Consejos concretos para estudiantes de doctorado:
\begin{itemize}
\item \textbf{Etapa uno} (los primeros 3 o 4 artículos): publicar los papers de manera honesta y disciplinada, «la etapa uno es competir al máximo en la ejecución». Obtener suficiente información y retroalimentación publicando muchos artículos.
\item \textbf{Etapa dos} (una vez que ya tienes una base): empezar a pensar en el why: «¿cómo salieron estos papers? ¿cuánta gente los cita? ¿forman un sistema coherente? ¿han trascendido su nicho?»; pasar de las publicaciones de relleno a hacer un trabajo con impact.
\end{itemize}

El criterio para juzgar si alguien puede llegar a ser un buen impact maker en la etapa dos: \textbf{ver si es capaz de reconocer que en el fondo es un pendejo: respetar los hechos, admitir el error cuando se equivoca, mejorar sin parar, y tener un buen role model, aunque ese role model esté muy lejos de él.}

\subsection{Explore vs. Exploit: visión de largo plazo y ejecución de corto plazo}

Hu Yuanming dedica cerca del 20\% de su tiempo a pensar en el futuro. Enfatiza que la vision no es algo que se hace una sola vez: «tal vez una palabra mejor sea visioning; tienes que estar pensando todo el tiempo en cómo es en realidad el futuro».

Su visión de futuro gira en torno a «AI for fun»:

\begin{knowledgebox}{AI for Fun: el producto definitivo de la era pos-AGI}
«Cuando llegue la AGI, todo el mundo va a perder su trabajo. Un mundo mejor es uno en el que la gente todavía pueda vivir feliz. Si con la tecnología de la AI generativa podemos crear para todos una immersive experience, y el tiempo limitado de cada persona se vuelve infinitamente rico gracias a la AI generativa —puedes volver al pasado, vivir vidas distintas—, creo que esto es, ultimately, una de las tecnologías más definitivas que el mundo necesita.»
\end{knowledgebox}

\subsection{La vida son diez grandes eventos}

Hu Yuanming tiene un marco singular para planear la vida: supongamos que empiezas a actuar a los 20 años y te retiras a los 70, con 50 años en medio. Si cada gran evento necesita 5 años, se pueden hacer 10 cosas. De esas, necesariamente la mitad va a fracasar: «si las 10 salen bien, eso significa que lo que hiciste tenía muy poco riesgo, y no era innovación de verdad».

Por ahora se pone una calificación de avance de «0.5 + 0.5»: haber llevado Taichi (太极) hasta donde está ahora vale 0.5, y haber llevado Meshy más allá del PMF vale otro 0.5. Sacar bien Taichi 2.0 es el siguiente 0.5, y llevar Meshy hasta la salida a bolsa es otro 0.5. Después de la salida a bolsa, ir por «el siguiente 1».

\subsection{Su manera de entender el fracaso}

\begin{importantbox}{Qué es el fracaso de verdad}
«¿Y qué si fracaso? ¿Me voy a morir? No. Mientras siga vivo, no he fracasado. Cualquier fracaso es solo algo temporal. \textbf{El fracaso de verdad} es dejar de querer innovar después de uno o dos intentos: que una serpiente te muerda una vez y durante diez años le tengas miedo hasta a una cuerda de pozo, y termines haciendo con disciplina solo cosas conservadoras. Si de 10 intentos de innovación 1 sale bien, ya es bastante bueno; fracasar es completamente normal.»
\end{importantbox}

\subsection{Consejos para los jóvenes}

\textbf{Para quienes quieren hacer un doctorado}: procurar encontrar en el camino algo que de verdad amen. «He visto a demasiada gente hacer un doctorado solo por el grado, y al final tampoco les gusta lo que hacen. Encontrar lo que amas no es algo que encuentres haciendo un doctorado: yo, antes de hacer el doctorado, ya había logrado decenas de artículos en SIGGRAPH, así que sabía que me iba a encantar.»

\textbf{Para quienes quieren emprender}:
\begin{enumerate}
\item \textbf{Get ready}: no se trata de tener claro qué vas a hacer, sino de tener la capacidad de soportar el fracaso: que tus padres estén sanos, que tu familia esté estable, poder pensar en la empresa las 24 horas.
\item \textbf{Take action}: la experiencia no importa tanto, muchas cosas se pueden aprender rápido por ensayo y error. «Si fracasas, fracasas, ¿y eso qué?»
\item Su propia mentalidad en ese momento era: «Si me gradúo del PhD en tres años y medio, aunque desperdicie dos años eso equivale a graduarme en cinco años y medio, ¿y eso qué?»
\end{enumerate}

\subsection{¿Todavía necesita un CEO escribir código?}

Hu Yuanming todavía dedica algunas horas por semana a escribir código (cerca del 10\% de su tiempo), sobre todo para ayudar al equipo a hacer debug y review del código. Citó lo que dijo Lisa Su, CEO de AMD: «su mayor fortaleza frente a un MBA es que actually understand cómo funciona la tecnología».

\begin{warningbox}{El peligro de que el CEO no escriba código}
«Si un CEO no escribe código, es muy fácil que se vuelva arrogante y crea que todo son problemas fáciles de resolver. Pero si de verdad te pones a ver cada día algunos de los "problemas de mierda" que la gente tiene que resolver, te das cuenta de que la cosa es bastante complicada. Un buen Founder siempre es hands-on: Lei Jun, Jensen Huang y Elon Musk son así. No necesariamente se trata de escribir código; puede ser diseñar el producto, hablar directamente con los proveedores, o tener un entendimiento profundo de la tecnología.»
\end{warningbox}

\subsection{Cápsula del tiempo: una palabra}

Al final de la entrevista, el conductor le pidió a Hu Yuanming que dejara una cápsula del tiempo para 2025. Él dijo una sola palabra: \textbf{valentía}.

«Espero que todos los que vean este Podcast tengan valentía.»

¿Qué haría si dejara de emprender? Dice que \textbf{escribiría un libro}: haría una «versión remasterizada en alta definición» del curso Games201, y escribiría un tutorial sobre un hybrid simulator de Simulation + Neural Network. «En su momento lo di demasiado apurado: todos los días de 8 a 9 de la mañana, y luego me iba a hacer mi práctica profesional en NVIDIA.»

\subsection{La distribución global de Meshy y la contratación de personal}

Meshy es una empresa globally distributed: el propio Hu Yuanming está en Silicon Valley, y el equipo está repartido entre Beijing, Shanghái, Shenzhen y otras partes del mundo. La empresa está pasando poco a poco del modelo work-from-home a un esquema híbrido: al menos 4 días a la semana en la office (en la office de China, los miércoles es teletrabajo), porque «discutir en persona es más eficiente».

En cuanto a la contratación, Meshy actualmente se enfoca en contratar AI Researcher, Machine Learning Engineer y Compiler Engineer (como reserva para Taichi 2.0). A los becarios se les paga con un salario global unificado (global pay), con un máximo de 1,000,000 de yuanes o 140,000 dólares al año, y el propio CEO participa directamente en guiar a los becarios en producto e investigación.

Esta manera de «que el CEO mismo guíe a los becarios» refleja la insistencia de Hu Yuanming en ser hands-on: él considera que esto es la ventaja competitiva central de una empresa con ADN técnico: \textbf{si el CEO no entiende los detalles técnicos, ¿por qué habría de trabajar para ti el resto de la gente?}

\subsection{Resumen de la sección}

Esta sección concentra el núcleo de la filosofía de vida de Hu Yuanming: la visión de talento en cuatro dimensiones Brain/Guts/Heart/Taste va más allá del marco tradicional de «inteligencia + esfuerzo»; el modelo de salto de la etapa uno a la etapa dos explica por qué muchos colaboradores individuales excelentes no logran ser buenos Managers; la redefinición del fracaso («mientras siga vivo, no he fracasado») da una base psicológica para la innovación constante. Y la palabra que le deja al mundo —valentía— es tanto una síntesis de su propio camino como una expectativa para quienes vienen después.
\section{Síntesis y ampliación}

\subsection{Mapa de ideas centrales}

La forma de pensar de Hu Yuanming (胡渊鸣) puede resumirse en unas cuantas convicciones centrales:

\textbf{Sobre la visión del mundo}: es muy probable que este mundo sea una simulación —la constante de Planck sería la precisión de punto flotante, y la velocidad de la luz, un límite de recursos de cómputo. Como integrantes de la simulación, no podemos refutar esta hipótesis desde dentro, pero esta perspectiva influye profundamente en su manera de entender la relación entre la simulación, la IA y el mundo real.

\textbf{Sobre la tecnología}: las leyes físicas son exactas, pero las condiciones de frontera no lo son. La simulación del futuro se apoyará sobre todo en el aprendizaje, con las leyes físicas como apoyo secundario. El 3D inductive bias es el puente que conecta la simulación con la red neuronal. No hay que subestimar la velocidad del avance tecnológico.

\textbf{Sobre el emprendimiento}: elegir la dirección correcta importa diez veces más que el esfuerzo. La tecnología solo tiene continuidad si puede convertirse en producto comercial. El código abierto puede ser un arma estratégica, pero debe complementarse con la conversión en producto. Las empresas emergentes deben buscar mercados de tamaño mediano y no consensuados.

\textbf{Sobre el crecimiento personal}: el 99\% es suerte (incluida la capacidad misma de esforzarse). Hay que dedicarse solo a lo que uno disfruta de verdad. Cada tres meses hay que renegar de quien uno era antes. Conviene dedicar más tiempo a pensar el qué y el por qué, no solo el cómo. No hay que aspirar a que todos te quieran.

\subsection{Lo que aporta a distintos lectores}

\textbf{Para quienes hacen investigación}: dedicar el 95\% del tiempo a pensar cuál es el problema importante, y solo el 5\% a resolverlo (cita de Einstein). Antes de publicar un artículo, hay que preguntarse «¿este trabajo importa?» —es también la pregunta que Hu Yuanming hace en cada entrevista de trabajo en Meshy—. No hay que encerrarse en la oficina: mantén la puerta abierta y escucha lo que hace la gente de otros campos. En la primera etapa hay que competir sin reservas cuando toca, pero después de tener de tres a cuatro artículos en conferencias de primer nivel, hay que cambiar a la mentalidad de la segunda etapa: no publicar más artículos, sino hacer un trabajo con más impacto.

\textbf{Para quienes emprenden}: la conversión en producto comercial es el primer principio —una tecnología que no puede convertirse en producto comercial es solo un «espectáculo de fuegos artificiales». El carácter determina el modelo de negocio —los fundadores introvertidos deben construir productos de consumo estandarizados (suscripción más API) y evitar la subcontratación, que exige mucho mantenimiento de relaciones con clientes. Un pivote no da miedo; lo que da miedo es no tener el valor de pivotar. Hay que atenerse a los hechos y respetar los datos, sin dejarse engañar por la sopa de pollo para el alma de que «el cielo recompensa al esforzado». Los tres elementos de un buen negocio: un mercado grande, no consensuado (cuestión de momento oportuno) y con ventaja competitiva. Elegir la dirección importa diez veces más que el esfuerzo.

\textbf{Para quienes dirigen equipos}: un buen líder no tiene por qué caerle bien a todos, pero debe ganarse el respeto de todos. El mayor obstáculo para quien asciende a un puesto directivo es atreverse a exigir estándares más altos a quienes antes eran sus pares. La cultura de la empresa está al servicio de «ganar» —el desgaste interno, pasarse la responsabilidad unos a otros y la búsqueda de armonía a toda costa son «culturas que llevan a perder». Lo más importante para un director general: 40\% en contratación, 30\% en pensar la estrategia futura y 30\% en participar en las operaciones diarias. Hay que mantenerse involucrado de manera directa —un director general que no escribe código tiende a «perder el piso».

\textbf{Para todos}:
\begin{itemize}
\item Encuentra lo que de verdad disfrutas, cuanto antes mejor
\item El 99\% es suerte: mantén la humildad
\item Revísate cada tres meses —si no te parece que quien eras hace tres meses era un idiota, es que no has crecido
\item No aspires a que todos te quieran: dedica tu energía a hacer lo correcto
\item El verdadero fracaso no es que un proyecto salga mal, sino no atreverse a intentarlo de nuevo
\item Hay que tener valor —«espero que todos los que vean este pódcast tengan valor»
\end{itemize}

\subsection{Selección de citas de Hu Yuanming}

A continuación, las frases originales más reveladoras de la entrevista:

\begin{table}[H]
\centering
\caption{Citas centrales de Hu Yuanming}
\begin{tabular}{p{3.5cm}p{10cm}}
\toprule
\textbf{Tema} & \textbf{Cita original} \\
\midrule
Autodefinición & «Soy una persona que se mueve por su propio interés y su sentido de misión.» \\
\midrule
Suerte & «El 99\% es suerte. Incluso el poder esforzarte es parte de la suerte.» \\
\midrule
Crecimiento en el emprendimiento & «Hay que despedirse a uno mismo cada tres o seis meses. Lo ideal es sentir que quien eras hace tres meses era un idiota.» \\
\midrule
Conciencia de director general & «Si el objetivo de tu trabajo es que todos te quieran, terminarás por caerle mal a todos.» \\
\midrule
Valor del código abierto & «El código abierto en efecto no genera ingresos, pero te ayuda a atraer al mejor talento y a los mejores clientes.» \\
\midrule
Esfuerzo y recompensa & «La inmensa mayoría del esfuerzo en el mundo no tiene ninguna recompensa. Eso de que el cielo recompensa al esforzado es pura mentira.» \\
\midrule
Visión del fracaso & «Mientras siga vivo, no he fracasado. El verdadero fracaso es no querer innovar más.» \\
\midrule
Forja del emprendimiento & «En tiempos de paz, nada forma tanto a una persona como emprender.» \\
\midrule
Hipótesis de la simulación & «Quizás el mundo entero sea un programa que una civilización más avanzada escribió por aburrimiento para simularnos.» \\
\midrule
IA y física & «El simulador necesita más IA, y la IA necesita más 3D inductive bias.» \\
\bottomrule
\end{tabular}
\end{table}

\subsection{Lecturas adicionales}

\begin{itemize}
\item \textbf{Lenguaje de programación Taichi}: \url{https://github.com/taichi-dev/taichi}—con más estrellas en GitHub que NVIDIA Warp, Triton y Mojo
\item \textbf{Meshy AI}: \url{https://www.meshy.ai/}—la plataforma de generación 3D con mayor participación de mercado actualmente
\item \textbf{Artículos de Hu Yuanming en Zhihu}: reflexiones sobre «si volviera a hacer el doctorado» y el planteamiento de «no ser solo un pez pequeño en SIGGRAPH»
\item \textbf{Andy Grove}: \textit{Solo los paranoicos sobreviven} (\textit{Only the Paranoid Survive})—el proceso de decisión de Intel al pasar de la DRAM a los CPU
\item \textbf{Richard Hamming}: \textit{You and Your Research}—la conferencia clásica sobre «mantener la puerta abierta»
\item \textbf{Sergey Levine}: \textit{The Spork of AGI}—sobre el papel y los límites de la simulación en la robótica
\item \textbf{Andrej Karpathy}: sobre la idea de que «un beneficio diez veces mayor suele costar solo el doble de esfuerzo»
\item \textbf{Games201}: el curso de simulación física que Hu Yuanming impartió en chino durante la pandemia de covid-19; planea lanzar en el futuro una «versión remasterizada en alta definición»
\item \textbf{DiffTaichi (ICLR 2020)}: el artículo emblemático de la programación diferenciable de Taichi, que une la simulación con la robótica
\item \textbf{ChainQueen}: el trabajo temprano de Taichi en simulación diferenciable
\item \textbf{QuantTaichi}: el trabajo de optimización por cuantización que logró simular mil millones de partículas en una sola GPU
\item \textbf{Lei Jun (雷军)}: \textit{Reflexiones sobre el emprendimiento de Xiaomi} (\textit{小米创业思考})—referencia de emprendimiento recomendada por Hu Yuanming, con la «teoría del cerdo volador» y la importancia de elegir la dirección
\end{itemize}

%% --- Fin del contenido --- %%

\end{document}
